problem:  
**Conjecture (Rigidity of locally homogeneous torsion-free Spin(7) structures).**  
Let \(M^8\) be a compact, locally homogeneous Riemannian manifold equipped with a torsion-free \(\mathrm{Spin}(7)\)-structure, so that the Levi–Civita connection satisfies  
\[
\nabla \Phi = 0, \qquad \mathrm{Hol}(M) = \mathrm{Spin}(7),
\]
where \(\Phi\) is the defining 4‑form of the \(\mathrm{Spin}(7)\) structure.  
Then \(M\) must be flat.  
Equivalently, there are **no nonflat, locally homogeneous Riemannian manifolds of dimension 8 with full holonomy \(\mathrm{Spin}(7)\).**

In a weaker form: classify all locally homogeneous Riemannian metrics \(g\) whose holonomy is contained in \(\mathrm{Spin}(7)\). Are all such metrics locally symmetric or flat?

---

Why is it a "good" problem:  
This conjecture isolates a strikingly simple yet unresolved rigidity question lying at the interface of **homogeneous geometry** and **special holonomy theory**.  

1. **Conceptual depth and simplicity.**  
   The statement—whether local homogeneity and exceptional holonomy are compatible—is immediately comprehensible to any differential geometer, yet touches the very core of the Berger classification and the Ambrose–Singer theory. It asks for a clean dichotomy: does the presence of parallel symmetry (local homogeneity) force curvature algebraic closure incompatible with Spin(7) reduction?

2. **Bridging distinct theories.**  
   Homogeneous spaces are controlled by **Lie algebraic curvature equations**, while torsion-free Spin(7) metrics are governed by **nonlinear PDEs on the defining 4-form**. The conjecture directly ties an analytic holonomy condition to an algebraic curvature condition, pushing toward a synthesis between **Lie group methods** and **special holonomy geometry**.

3. **Relation to known results but open direction.**  
   While existing works (e.g., by Fernández, Fino, Agricola, Nagy) classify *invariant* Spin(7)-structures on specific Lie groups and show that none are torsion-free except the flat ones, a **general theorem** banning *all* locally homogeneous, but not necessarily globally homogeneous, torsion-free Spin(7) geometries does not exist. Hence, extending such obstructions from global group settings to abstract local homogeneity remains genuinely unresolved and would represent a conceptual unification of current partial results.

4. **Potential implications.**  
   A proof would yield a powerful **rigidity theorem**: exceptional holonomy metrics cannot even locally admit a transitive Lie algebra of Killing fields except in the flat case—establishing a clear geometric divide between “algebraically symmetric” and “analytically exceptional” Riemannian geometries. Conversely, a counterexample would be revolutionary, identifying an unprecedented algebraic model with Spin(7) holonomy.

Thus, the conjecture exemplifies the ideal of a good research problem: conceptually simple, deeply linked to major geometric frameworks, unifying algebraic and analytic structures, and of plausibly open status at the frontier of differential geometry.